\documentclass[12pt]{article}
\usepackage[T1]{fontenc}
\usepackage[utf8]{inputenc}
\usepackage{lmodern}
\usepackage{textcomp}
\usepackage{enumitem}
\usepackage{geometry}
\geometry{margin=1in}
\begin{document}
\begin{center}
Answer Key - Constitutional Law - Undergraduate Level Assessment 24 \
\small (Answers are ordered exactly as in the question paper.)
\end{center}

\section*{Multiple Choice}
\begin{enumerate}[label=\arabic*.]
\item \textbf{Answer:} B
\\ \textit{Options:} A) Recognition of governments is very prevalent in contemporary practice; B) Recognition of governments has largely been replaced by functional Recognition; C) Government recognition is common in respect of rebel entities; D) Only democratic governments are recognised in contemporary practice
\\ \textit{Note:} Recognition of governments has largely been replaced by functional recognition because contemporary international practice focuses on the legitimacy of the governing authority and its capacity to fulfill international obligations, rather than merely recognizing the government itself.
\item \textbf{Answer:} C
\\ \textit{Options:} A) Cause and effect; B) Theory and fact; C) Being obliged and having an obligation; D) Corporeal and incorporeal rights
\\ \textit{Note:} Hart's analysis of law distinguishes between being obliged, which refers to the actual behavior compelled by external factors, and having an obligation, which reflects the internal moral or legal duty recognized by the individual, highlighting the difference between coercion and normative expectations in legal theory.
\item \textbf{Answer:} B
\\ \textit{Options:} A) ordinances; B) federal statutes; C) executive orders; D) charters
\\ \textit{Note:} The U.S. Congress is granted the authority to enact federal statutes under the Commerce Clause of the Constitution, which allows it to regulate activities that affect foreign and interstate commerce.
\item \textbf{Answer:} C
\\ \textit{Options:} A) Command, sovereign and enforceability; B) Command, sovereign and legal remedy; C) Command, sovereign and sanction; D) Command, sovereign and obedience by subject
\\ \textit{Note:} The correct answer identifies the key attributes of law as defined by legal philosopher John Austin, who emphasized that law is a command issued by a sovereign authority and backed by the threat of sanction for non-compliance.
\item \textbf{Answer:} D
\\ \textit{Options:} A) Goodwin v UK (2002); B) Airey v Ireland (1979); C) Osman v UK ( 1998); D) Handyside v UK (1976)
\\ \textit{Note:} The correct answer is Handyside v UK (1976) because this landmark case established the doctrine of 'margin of appreciation,' allowing states a degree of discretion in balancing individual rights with public interests under the European Convention on Human Rights.
\end{enumerate}

\section*{True / False}
\begin{enumerate}[label=\arabic*.]
\setcounter{enumi}{5}
\item \textbf{Answer:} True
\\ \textit{Options:} A) True; B) False
\\ \textit{Note:} The rule of recognition establishes the criteria for identifying valid legal norms within a legal system, thereby granting judges the authority to interpret and apply the law while also obligating them to resolve disputes in accordance with those recognized norms.
\item \textbf{Answer:} False
\\ \textit{Options:} A) True; B) False
\\ \textit{Note:} The statement is false because laws can and often do incorporate ethical considerations, reflecting societal values and moral principles alongside legal formalities.
\item \textbf{Answer:} False
\\ \textit{Options:} A) True; B) False
\\ \textit{Note:} Hart's distinction between law and morality argues that while they are separate, they can still influence each other, thus not entirely refuting the idea of their connection in legal systems.
\item \textbf{Answer:} True
\\ \textit{Options:} A) True; B) False
\\ \textit{Note:} The statement is true because the principle of sovereign immunity, rooted in constitutional law, protects states from being sued in the courts of other states, thereby upholding their sovereignty and dignity within the federal system.
\item \textbf{Answer:} True
\\ \textit{Options:} A) True; B) False
\\ \textit{Note:} The Social Fact Thesis is true because it posits that laws are not inherent truths but rather products of social agreements and cultural practices, reflecting the collective values and norms of society.
\end{enumerate}

\section*{Long Form Response}
\begin{enumerate}[label=\arabic*.]
\setcounter{enumi}{10}
\item \textbf{Answer:} Hobbes' assertion that peace is the first law of nature underscores the fundamental necessity of societal order for the protection of individual rights. In a state of nature, where peace is absent, individuals possess an unrestrained right to everything, including the lives of others, leading to a perpetual state of conflict and insecurity. This chaotic condition necessitates the establishment of a social contract, wherein individuals consent to surrender certain freedoms in exchange for security and order, thereby enabling the protection of their remaining rights. Within the framework of constitutional law, this concept emphasizes the importance of creating a governing system that prioritizes peace as a means to safeguard individual liberties and maintain societal stability. Ultimately, Hobbes' perspective highlights the delicate balance between individual rights and the collective need for order in any constitutional framework.
\\ \textit{Note:} The answer correctly identifies that Hobbes views peace as essential for societal order, arguing that without it, individual rights are jeopardized, necessitating a social contract to ensure security and stability within a constitutional framework.
\item \textbf{Answer:} Maine's aphorism that 'the movement of progressive societies has hitherto been a movement from Status to Contract' highlights a historical transition in social organization, emphasizing the shift from rigid hierarchies to more flexible, voluntary agreements. However, its misinterpretation as a prediction suggests that societies are destined to continue on this trajectory indefinitely, which oversimplifies the complexities of societal evolution. This misunderstanding can lead to a deterministic view of progress, ignoring the potential for regression or the emergence of new forms of status-based systems. Therefore, recognizing this aphorism as a reflection of past trends rather than a forecast allows for a more nuanced understanding of societal dynamics and the factors that influence change in constitutional law.
\\ \textit{Note:} The answer correctly identifies that Maine's aphorism illustrates a shift from rigid social structures to more flexible agreements, while also critiquing the flawed notion that this evolution is a linear and inevitable process, thereby emphasizing the complexities of societal change.
\end{enumerate}

\vfill
\noindent\textit{End of Answer Key}
\end{document}